%% kapa: exam=1 round=24 year=2013 period=1 form=A subject=economics subject_ko=경제원론 seq=12 number=52
%% source: Q-Net 감정평가사 기출문제 (cst00309) artlSeq=5073802 file=extracted/1008/5073802/2013년 제24회 감정평가사 1차시험 기출문제/1교시(A형).hwp
%% conversion: kordoc 4.24.51 -> tools/gen1.py (auto); review=proofread-sonnet; flags=none
\begin{question}{52}
독점시장에 관한 설명으로 옳지 않은 것은? [단, 수요곡선은 우하향하고 생산량은 양(+)이다.]
\begin{choices}
  \item 독점기업은 시장지배력을 이용하여 가격을 인상하면서 판매량을 늘릴 수 있다.
  \item 수요의 가격탄력성이 1일 때 독점기업의 한계수입은 0이다.
  \item 생산량이 증가함에 따라 평균비용곡선이 지속적으로 우하향하는 산업에서는 자연독점이 발생한다.
  \item 독점기업의 한계수입곡선은 수요곡선보다 아래에 위치한다.
  \item 독점기업의 이윤극대화 가격은 한계비용보다 높다.
\end{choices}
\end{question}
\answer{1}
